\ifdefined\gkver\else\def\gkver{student}\fi
\documentclass[\gkver]{gaokaozhenti}
\begin{document}

\chapter{2022年全国甲卷}
%% paperType: 理综物理部分

\begin{enumerate}
\item
%% number: 1
%% paperName: 2022年全国高考甲卷第 1 题 6 分
%% typeId: 1
%% type: 单选题
%% chapter: 曲线运动
%% point: 竖直平面圆周运动
%% method: 无
%% score: 6
%% degree: 750
%% duplicateId: 0
%% body:
北京2022年冬奥会首钢滑雪大跳台局部示意图如图所示。运动员从 a 处由静止自由滑下，到 b 处起跳，c 点为 a、b 之间的最低点，a、c 两处的高度差为 $h$。要求运动员经过 c 点时对滑雪板的压力不大于自身所受重力的 $k$ 倍，运动过程中将运动员视为质点并忽略所有阻力，则 c 点处这一段圆弧雪道的半径不应小于\xzanswer{D}

\onepicture[w=7cm]{figs/01.png}

\fourchoices[answer=D]
{$\frac{h}{k+1}$}
{$\frac{h}{k}$}
{$\frac{2h}{k}$}
{$\frac{2h}{k-1}$}

%% answer: D
%% memo:
\memoanswer{
运动员从 a 到 c 根据动能定理有

$mgh=\frac{1}{2}mv_{c}^{2}$

在 c 点有

$F_{\mathrm{Nc}}-mg=m\frac{v_{c}^{2}}{R_{c}}$

$F_{\mathrm{Nc}}\leq kmg$

联立有：$R_{c}\geqslant\frac{2h}{k-1}$

故选 D。
}

\item
%% number: 2
%% paperName: 2022年全国高考甲卷第 2 题 6 分
%% typeId: 1
%% type: 单选题
%% chapter: 直线运动
%% point: 多段匀变速直线运动
%% method: 无
%% score: 6
%% degree: 550
%% duplicateId: 0
%% body:
长为 $l$ 的高速列车在平直轨道上正常行驶，速率为 $v_{0}$，要通过前方一长为 $L$ 的隧道，当列车的任一部分处于隧道内时，列车速率都不允许超过 $v$（$v<v_{0}$）。已知列车加速和减速时加速度的大小分别为 $a$ 和 $2a$，则列车从减速开始至回到正常行驶速率 $v_{0}$ 所用时间至少为\xzanswer{C}

\fourchoices[answer=C]
{$\frac{v_{0}-v}{2a}+\frac{L+l}{v}$}
{$\frac{v_{0}-v}{2a}+\frac{L+2l}{v}$}
{$\frac{3(v_{0}-v)}{2a}+\frac{L+l}{v}$}
{$\frac{3(v_{0}-v)}{2a}+\frac{L+2l}{v}$}

%% answer: C
%% memo:
\memoanswer{
由题知当列车的任一部分处于隧道内时，列车速率都不允许超过 $v$（$v<v_{0}$），则列车进隧道前必须减速到 $v$，则有

$v=v_{0}-2at_{1}$

解得

$t_{1}=\frac{v_{0}-v}{2a}$

在隧道内匀速有

$t_{2}=\frac{L+l}{v}$

列车尾部出隧道后立即加速到 $v_{0}$，有

$v_{0}=v+at_{3}$

解得

$t_{3}=\frac{v_{0}-v}{a}$

则列车从减速开始至回到正常行驶速率 $v_{0}$ 所用时间至少为

$t=\frac{3(v_{0}-v)}{2a}+\frac{L+l}{v}$

故选 C。
}

\item
%% number: 3
%% paperName: 2022年全国高考甲卷第 3 题 6 分
%% typeId: 1
%% type: 单选题
%% chapter: 电磁感应
%% point: 法拉第电磁感应定律
%% method: 无
%% score: 6
%% degree: 550
%% duplicateId: 0
%% body:
三个用同样的细导线做成的刚性闭合线框，正方形线框的边长与圆线框的直径相等，圆线框的半径与正六边形线框的边长相等，如图所示。把它们放入磁感应强度随时间线性变化的同一匀强磁场中，线框所在平面均与磁场方向垂直，正方形、圆形和正六边形线框中感应电流的大小分别为 $I_{1}$、$I_{2}$ 和 $I_{3}$。则\xzanswer{C}

\onepicture[w=8cm]{figs/03.png}

\fourchoices[answer=C]
{$I_{1}<I_{3}<I_{2}$}
{$I_{1}>I_{3}>I_{2}$}
{$I_{1}=I_{2}>I_{3}$}
{$I_{1}=I_{2}=I_{3}$}

%% answer: C
%% memo:
\memoanswer{
设圆线框的半径为 $r$，则由题意可知正方形线框的边长为 $2r$，正六边形线框的边长为 $r$；所以圆线框的周长为 $C_{2}=2\pi r$，面积为 $S_{2}=\pi r^{2}$；同理可知正方形线框的周长和面积分别为

$C_{1}=8r$，$S_{1}=4r^{2}$

正六边形线框的周长和面积分别为

$C_{3}=6r$，$S_{3}=\frac{1}{2}\times6\times r\times\frac{\sqrt{3}}{2}r=\frac{3\sqrt{3}r^{2}}{2}$

三线框材料粗细相同，根据电阻定律

$R=\rho\frac{L}{S_{\text{横截面}}}$

可知三个线框电阻之比为

$R_{1}:R_{2}:R_{3}=C_{1}:C_{2}:C_{3}=8:2\pi:6$

根据法拉第电磁感应定律有

$I=\frac{E}{R}=\frac{\Delta B}{\Delta t}\cdot\frac{S}{R}$

可得电流之比为：

$I_{1}:I_{2}:I_{3}=2:2:\sqrt{3}$

即

$I_{1}=I_{2}>I_{3}$

故选 C。
}

\item
%% number: 4
%% paperName: 2022年全国高考甲卷第 4 题 6 分
%% typeId: 1
%% type: 单选题
%% chapter: 原子和原子核
%% point: 半衰期
%% method: 无
%% score: 6
%% degree: 450
%% duplicateId: 0
%% body:
两种放射性元素的半衰期分别为 $t_{0}$ 和 $2t_{0}$，在 $t=0$ 时刻这两种元素的原子核总数为 $N$，在 $t=2t_{0}$ 时刻，尚未衰变的原子核总数为 $\frac{N}{3}$，则在 $t=4t_{0}$ 时刻，尚未衰变的原子核总数为\xzanswer{C}

\fourchoices[answer=C]
{$\frac{N}{12}$}
{$\frac{N}{9}$}
{$\frac{N}{8}$}
{$\frac{N}{6}$}

%% answer: C
%% memo:
\memoanswer{
根据题意设半衰期为 $t_{0}$ 的元素原子核数为 $x$，另一种元素原子核数为 $y$，依题意有

$x+y=N$

经历 $2t_{0}$ 后有

$\frac{1}{4}x+\frac{1}{2}y=\frac{N}{3}$

联立可得

$x=\frac{2}{3}N$，$y=\frac{1}{3}N$

在 $t=4t_{0}$ 时，原子核数为 $x$ 的元素经历了 4 个半衰期，原子核数为 $y$ 的元素经历了 2 个半衰期，则此时未衰变的原子核总数为

$n=\frac{1}{2^{4}}x+\frac{1}{2^{2}}y=\frac{N}{8}$

故选 C。
}

\item
%% number: 5
%% paperName: 2022年全国高考甲卷第 5 题 6 分
%% typeId: 1
%% type: 单选题
%% chapter: 磁场
%% point: 带电粒子在复合场中的运动
%% method: 无
%% score: 6
%% degree: 450
%% duplicateId: 0
%% body:
空间存在着匀强磁场和匀强电场，磁场的方向垂直于纸面（$xOy$ 平面）向里，电场的方向沿 $y$ 轴正方向。一带正电的粒子在电场和磁场的作用下，从坐标原点 $O$ 由静止开始运动。下列四幅图中，可能正确描述该粒子运动轨迹的是\xzanswer{B}

\fourchoices[answer=B, ispicture=true, h=2.3cm]
{figs/05a.jpg}
{figs/05b.jpg}
{figs/05c.jpg}
{figs/05d.jpg}

%% answer: B
%% memo:
\memoanswer{
解法一：

AC．在 $xOy$ 平面内电场的方向沿 $y$ 轴正方向，故在坐标原点 $O$ 静止的带正电粒子在电场力作用下会向 $y$ 轴正方向运动。磁场方向垂直于纸面向里，根据左手定则，可判断出向 $y$ 轴正方向运动的粒子同时受到沿 $x$ 轴负方向的洛伦兹力，故带电粒子向 $x$ 轴负方向偏转。AC 错误；

BD．运动的过程中在电场力对带电粒子做功，粒子速度大小发生变化，粒子所受的洛伦兹力方向始终与速度方向垂直。由于匀强电场方向是沿 $y$ 轴正方向，故 $x$ 轴为匀强电场的等势面，从开始到带电粒子偏转再次运动到 $x$ 轴时，电场力做功为 0，洛伦兹力不做功，故带电粒子再次回到 $x$ 轴时的速度为 0，随后受电场力作用再次进入第二象限重复向左偏转，故 B 正确，D 错误。

故选 B。

解法二：

粒子在 $O$ 点静止，对速度进行分解，分解为向 $x$ 轴正方向的速度 $v$，向 $x$ 轴负方向的速度 $v'$，两个速度大小相等，方向相反。使得其中一个洛伦兹力平衡电场力，即 $qBv'=qE$

则粒子在电场、磁场中的运动，可视为，向 $x$ 轴负方向以速度 $v'=\frac{E}{B}$ 做匀速直线运动，同时在 $x$ 轴上方做匀速圆周运动。

故选 B。
}

\item
%% number: 6
%% paperName: 2022年全国高考甲卷第 6 题 6 分
%% typeId: 2
%% type: 多选题
%% chapter: 运动和力
%% point: 变加速直线运动
%% method: 无
%% score: 6
%% degree: 450
%% duplicateId: 0
%% body:
如图，质量相等的两滑块 P、Q 置于水平桌面上，二者用一轻弹簧水平连接，两滑块与桌面间的动摩擦因数均为 $\mu$。重力加速度大小为 $g$。用水平向右的拉力 $F$ 拉动 P，使两滑块均做匀速运动；某时刻突然撤去该拉力，则从此刻开始到弹簧第一次恢复原长之前\xzanswer{AD}

\onepicture[w=9cm]{figs/06.png}

\fourchoices[answer=AD]
{P 的加速度大小的最大值为 $2\mu g$}
{Q 的加速度大小的最大值为 $2\mu g$}
{P 的位移大小一定大于 Q 的位移大小}
{P 的速度大小均不大于同一时刻 Q 的速度大小}

%% answer: AD
%% memo:
\memoanswer{
设两物块的质量均为 $m$，撤去拉力前，两滑块均做匀速直线运动，则拉力大小为 $F=2\mu mg$

撤去拉力前对 Q 受力分析可知，弹簧的弹力为

$T_{0}=\mu mg$

AB．以向右为正方向，撤去拉力瞬间弹簧弹力不变为 $\mu mg$，两滑块与地面间仍然保持相对滑动，此时滑块 P 的加速度为

$-T_{0}-\mu mg=ma_{P1}$

解得

$a_{P1}=-2\mu g$

此刻滑块 Q 所受的外力不变，加速度仍为零，滑块 P 做减速运动，故 PQ 间距离减小，弹簧的伸长量变小，弹簧弹力变小。根据牛顿第二定律可知 P 减速的加速度减小，滑块 Q 的合外力增大，合力向左，做加速度增大的减速运动。

故 P 加速度大小的最大值是刚撤去拉力瞬间的加速度为 $2\mu g$。

Q 加速度大小最大值为弹簧恢复原长时

$-\mu mg=ma_{Qm}$

解得

$a_{Qm}=-\mu g$

故滑块 Q 加速度大小最大值为 $\mu g$，A 正确，B 错误；

C．滑块 PQ 水平向右运动，PQ 间的距离在减小，故 P 的位移一定小于 Q 的位移，C 错误；

D．滑块 P 在弹簧恢复到原长时的加速度为

$-\mu mg=ma_{P2}$

解得

$a_{P2}=-\mu g$

撤去拉力时，PQ 的初速度相等，滑块 P 由开始的加速度大小为 $2\mu g$ 做加速度减小的减速运动，最后弹簧原长时加速度大小为 $\mu g$；滑块 Q 由开始的加速度为 0 做加速度增大的减速运动，最后弹簧原长时加速度大小也为 $\mu g$。分析可知 P 的速度大小均不大于同一时刻 Q 的速度大小，D 正确。

故选 AD。
}

\item
%% number: 7
%% paperName: 2022年全国高考甲卷第 7 题 6 分
%% typeId: 2
%% type: 多选题
%% chapter: 电磁感应
%% point: 电磁感应含容电路
%% method: 无
%% score: 6
%% degree: 350
%% duplicateId: 0
%% body:
如图，两根相互平行的光滑长直金属导轨固定在水平绝缘桌面上，在导轨的左端接入电容为 $C$ 的电容器和阻值为 $R$ 的电阻。质量为 $m$、阻值也为 $R$ 的导体棒 MN 静止于导轨上，与导轨垂直，且接触良好，导轨电阻忽略不计，整个系统处于方向竖直向下的匀强磁场中。开始时，电容器所带的电荷量为 $Q$，合上开关 S 后，\xzanswer{AD}

\onepicture[w=7cm]{figs/07.png}

\fourchoices[answer=AD]
{通过导体棒 MN 电流的最大值为 $\frac{Q}{RC}$}
{导体棒 MN 向右先加速、后匀速运动}
{导体棒 MN 速度最大时所受的安培力也最大}
{电阻 $R$ 上产生的焦耳热大于导体棒 MN 上产生的焦耳热}

%% answer: AD
%% memo:
\memoanswer{
MN 在运动过程中为非纯电阻，MN 上的电流瞬时值为

$i=\frac{u-Blv}{R}$

A．当闭合的瞬间，$Blv=0$，此时 MN 可视为纯电阻 R，此时反电动势最小，故电流最大

$I_{max}=\frac{U}{R}=\frac{Q}{CR}$

故 A 正确；

B．当 $u>Blv$ 时，导体棒加速运动，当速度达到最大值之后，电容器与 MN 及 R 构成回路，由于一直处于通路的形式，由能量守恒可知，最后 MN 终极速度为零，故 B 错误；

C．MN 在运动过程中为非纯电阻电路，MN 上的电流瞬时值为

$i=\frac{u-Blv}{R}$

当 $u=Blv$ 时，MN 上电流瞬时为零，安培力为零此时，MN 速度最大，故 C 错误；

D．在 MN 加速度阶段，由于 MN 反电动势存在，故 MN 上电流小于电阻 R 上的电流，电阻 R 消耗电能大于 MN 上消耗的电能（即 $E_{R}>E_{MN}$），故加速过程中，$Q_{R}>Q_{MN}$；当 MN 减速为零的过程中，电容器的电流和导体棒的电流都流经电阻 R 形成各自的回路，因此可知此时也是电阻 R 的电流大，综上分析可知全过程中电阻 R 上的热量大于导体棒上的热量，故 D 正确。

故选 AD。
}

\item
%% number: 8
%% paperName: 2022年全国高考甲卷第 8 题 6 分
%% typeId: 2
%% type: 多选题
%% chapter: 电场
%% point: 带电粒子的偏转
%% method: 等效替代
%% score: 6
%% degree: 350
%% duplicateId: 0
%% body:
地面上方某区域存在方向水平向右的匀强电场，将一带正电荷的小球自电场中 P 点水平向左射出。小球所受的重力和电场力的大小相等，重力势能和电势能的零点均取在 P 点。则射出后，\xzanswer{BD}

\fourchoices[answer=BD]
{小球的动能最小时，其电势能最大}
{小球的动能等于初始动能时，其电势能最大}
{小球速度的水平分量和竖直分量大小相等时，其动能最大}
{从射出时刻到小球速度的水平分量为零时，重力做的功等于小球电势能的增加量}

%% answer: BD
%% memo:
\memoanswer{
A．如图所示

\onepicture[w=5.5cm]{figs/08_an.png}

$qE=mg$

故等效重力 $G'$ 的方向与水平成 $45^{\circ}$。

当 $v_{y}=0$ 时速度最小为 $v_{\min}=v_{1}$，由于此时 $v_{1}$ 存在水平分量，电场力还可以向左做负功，故此时电势能不是最大，故 A 错误；

BD．水平方向上：$v_{0}=\frac{qE}{m}t$

在竖直方向上：$v=gt$

由于 $qE=mg$，得 $v=v_{0}$

如图所示，小球的动能等于末动能。由于此时速度没有水平分量，故电势能最大。由动能定理可知

$W_{G}+W_{\text{电}}=0$

则重力做功等于小球电势能的增加量，故 BD 正确；

C．当如图中 $v_{1}$ 所示时，此时速度水平分量与竖直分量相等，动能最小，故 C 错误；

故选 BD。
}

\item
%% number: 9
%% paperName: 2022年全国高考甲卷第 9 题 5 分
%% typeId: 4
%% type: 实验题
%% chapter: 电路
%% point: 伏安法测电阻
%% method: 无
%% score: 5
%% degree: 650
%% duplicateId: 0
%% body:
某同学要测量微安表内阻，可利用的实验器材有：电源 $E$（电动势 $1.5\UV$，内阻很小），电流表（量程 $10\UmA$，内阻约 $10\UO$），微安表（量程 $100\UuA$，内阻 $R_{g}$ 待测，约 $1\UkO$），滑动变阻器 $R$（最大阻值 $10\UO$），定值电阻 $R_{0}$（阻值 $10\UO$），开关 S，导线若干。
\begin{enumerate}
	\item
	在答题卡上将图中所示的器材符号连线，画出实验电路原理图\tkanswer[2cm]{}；
	\item
	某次测量中，微安表的示数为 $90.0\UuA$，电流表的示数为 $9.00\UmA$，由此计算出微安表内阻 $R_{g}=$\tkanswer[1.8cm]{$990$}$\UO$。
\end{enumerate}

\onepicture[align=right, w=5.5cm]{figs/09.png}

%% answer: 
\jdanswer{
\begin{enumerate}
	\item
	如图

	\item
	$990$
\end{enumerate}
}
%% memo:
\memoanswer{
（1）为了准确测出微安表两端的电压，可以让微安表与定值电阻 $R_{0}$ 并联，再与电流表串联，通过电流表的电流与微安表的电流之差，可求出流过定值电阻 $R_{0}$ 的电流，从而求出微安表两端的电压，进而求出微安表的内电阻，由于电源电压过大，并且为了测量多组数据，滑动变阻器采用分压式解法，实验电路原理图如图所示

\onepicture[w=5cm]{figs/09_an.png}

（2）流过定值电阻 $R_{0}$ 的电流

$I=I_{A}-I_{G}=9.00\UmA-0.09\UmA=8.91\UmA$

加在微安表两端的电压

$U=IR_{0}=8.91\times10^{-2}\UV$

微安表的内电阻

$R_{g}=\frac{U}{I_{G}}=\frac{8.91\times10^{-2}}{90.0\times10^{-6}}\UO=990\UO$
}

\item
%% number: 10
%% paperName: 2022年全国高考甲卷第 10 题 10 分
%% typeId: 4
%% type: 实验题
%% chapter: 运动和力
%% point: 验证动量守恒实验
%% method: 无
%% score: 10
%% degree: 450
%% duplicateId: 0
%% body:
利用图示的实验装置对碰撞过程进行研究。让质量为 $m_{1}$ 的滑块 A 与质量为 $m_{2}$ 的静止滑块 B 在水平气垫导轨上发生碰撞，碰撞时间极短，比较碰撞后 A 和 B 的速度大小 $v_{1}$ 和 $v_{2}$，进而分析碰撞过程是否为弹性碰撞。完成下列填空：

\onepicture[align=right, w=11cm]{figs/10.png}

\begin{enumerate}
	\item
	调节导轨水平；
	\item
	测得两滑块的质量分别为 $0.510\Ukg$ 和 $0.304\Ukg$。要使碰撞后两滑块运动方向相反，应选取质量为\tkanswer[1.8cm]{0.304}$\Ukg$的滑块作为 A；
	\item
	调节 B 的位置，使得 A 与 B 接触时，A 的左端到左边挡板的距离 $s_{1}$ 与 B 的右端到右边挡板的距离 $s_{2}$ 相等；
	\item
	使 A 以一定的初速度沿气垫导轨运动，并与 B 碰撞，分别用传感器记录 A 和 B 从碰撞时刻开始到各自撞到挡板所用的时间 $t_{1}$ 和 $t_{2}$；
	\item
	将 B 放回到碰撞前的位置，改变 A 的初速度大小，重复步骤（4）。多次测量的结果如下表所示；
	\begin{table}[htbp]
	\centering
	\begin{tblr}{|c|c|c|c|c|c|}
	\hline
	 & 1 & 2 & 3 & 4 & 5 \\
	\hline
	$t_{1}/\Us[a]$ & 0.49 & 0.67 & 1.01 & 1.22 & 1.39 \\
	\hline
	$t_{2}/\Us[a]$ & 0.15 & 0.21 & 0.33 & 0.40 & 0.46 \\
	\hline
	$k_{2}=\frac{v_{1}}{v_{2}}$ & 0.31 & $k_{2}$ & 0.33 & 0.33 & 0.33 \\
	\hline
	\end{tblr}
	\end{table}
	\item
	表中的\tkanswer[1.8cm]{0.31}（保留 2 位有效数字）；
	\item
	$\frac{v_{1}}{v_{2}}$ 的平均值为\tkanswer[1.8cm]{0.32}；（保留 2 位有效数字）
	\item
	理论研究表明，对本实验的碰撞过程，是否为弹性碰撞可由 $\frac{v_{1}}{v_{2}}$ 判断。若两滑块的碰撞为弹性碰撞，则 $\frac{v_{1}}{v_{2}}$ 的理论表达式为\tkanswer[2.6cm]{$\frac{m_{2}-m_{1}}{2m_{1}}$}（用 $m_{1}$ 和 $m_{2}$ 表示），本实验中其值为\tkanswer[1.8cm]{0.34}（保留 2 位有效数字），若该值与（7）中结果间的差别在允许范围内，则可认为滑块 A 与滑块 B 在导轨上的碰撞为弹性碰撞。
\end{enumerate}

%% answer: 
\jdanswer{
\begin{enumerate}
	\item
	（无作答要求）

	\item
	$0.304$

	\item
	（无作答要求）

	\item
	（无作答要求）

	\item
	（无作答要求）

	\item
	$0.31$

	\item
	$0.32$

	\item
	$\frac{m_{2}-m_{1}}{2m_{1}}$，$0.34$
\end{enumerate}
}
%% memo:
\memoanswer{
（2）应该用质量较小的滑块碰撞质量较大的滑块，碰后运动方向相反，故选 $0.304\Ukg$ 的滑块作为 A。

（6）由于两段位移大小相等，根据表中的数据可得

$k_{2}=\frac{v_{1}}{v_{2}}=\frac{t_{2}}{t_{1}}=\frac{0.21}{0.67}=0.31$

（7）$\frac{v_{1}}{v_{2}}$ 平均值为

$\overline{k}=\frac{0.31+0.31+0.33+0.33+0.33}{5}=0.32$

（8）弹性碰撞时满足动量守恒和机械能守恒，可得

$m_{1}v_{0}=-m_{1}v_{1}+m_{2}v_{2}$

$\frac{1}{2}m_{1}v_{0}^{2}=\frac{1}{2}m_{1}v_{1}^{2}+\frac{1}{2}m_{2}v_{2}^{2}$

联立解得

$\frac{v_{1}}{v_{2}}=\frac{m_{2}-m_{1}}{2m_{1}}$

代入数据可得

$\frac{v_{1}}{v_{2}}=0.34$
}

\item
%% number: 11
%% paperName: 2022年全国高考甲卷第 11 题 12 分
%% typeId: 6
%% type: 计算题
%% chapter: 曲线运动
%% point: 平抛运动
%% method: 无
%% score: 12
%% degree: 450
%% duplicateId: 0
%% body:
将一小球水平抛出，使用频闪仪和照相机对运动的小球进行拍摄，频闪仪每隔 $0.05\Us$ 发出一次闪光。某次拍摄时，小球在抛出瞬间频闪仪恰好闪光，拍摄的照片编辑后如图所示。图中的第一个小球为抛出瞬间的影像，每相邻两个球之间被删去了 3 个影像，所标出的两个线段的长度 $s_{1}$ 和 $s_{2}$ 之比为 $3:7$。重力加速度大小取 $g=10\Umsq$，忽略空气阻力。求在抛出瞬间小球速度的大小。

\onepicture[align=right, w=6cm]{figs/11.png}

%% answer: 
\jdanswer{
$\frac{2\sqrt 5}{5}\Ums$
}
%% memo:
\memoanswer{
频闪仪每隔 $0.05\Us$ 发出一次闪光，每相邻两个球之间被删去 3 个影像，故相邻两球的时间间隔为 $t=4T=0.05\times4\Us=0.2\Us$

设抛出瞬间小球的速度为 $v_{0}$，每相邻两球间的水平方向上位移为 $x$，竖直方向上的位移分别为 $y_{1}$、$y_{2}$，根据平抛运动位移公式有

$x=v_{0}t$

$y_{1}=\frac{1}{2}gt^{2}=\frac{1}{2}\times10\times0.2^{2}\Um=0.2\Um$

$y_{2}=\frac{1}{2}g(2t)^{2}-\frac{1}{2}gt^{2}=\frac{1}{2}\times10\times(0.4^{2}-0.2^{2})\Um=0.6\Um$

令 $y_{1}=y$，则有

$y_{2}=3y_{1}=3y$

已标注的线段 $s_{1}$、$s_{2}$ 分别为

$s_{1}=\sqrt{x^{2}+y^{2}}$

$s_{2}=\sqrt{x^{2}+(3y)^{2}}=\sqrt{x^{2}+9y^{2}}$

则有

$\sqrt{x^{2}+y^{2}}:\sqrt{x^{2}+9y^{2}}=3:7$

整理得

$x=\frac{2\sqrt 5}{5}y$

故在抛出瞬间小球的速度大小为

$v_{0}=\frac{x}{t}=\frac{2\sqrt 5}{5}\Ums$
}

\item
%% number: 12
%% paperName: 2022年全国高考甲卷第 12 题 20 分
%% typeId: 6
%% type: 计算题
%% chapter: 磁场
%% point: 安培力
%% method: 无
%% score: 20
%% degree: 350
%% duplicateId: 0
%% body:
光点式检流计是一种可以测量微小电流的仪器，其简化的工作原理示意图如图所示。图中 A 为轻质绝缘弹簧，C 为位于纸面上的线圈，虚线框内有与纸面垂直的匀强磁场；M 为置于平台上的轻质小平面反射镜，轻质刚性细杆 D 的一端与 M 固连且与镜面垂直、另一端与弹簧下端相连，PQ 为圆弧形的、带有均匀刻度的透明读数条，PQ 的圆心位于 M 的中心。使用前需调零：使线圈内没有电流通过时，M 竖直且与纸面垂直；入射细光束沿水平方向经 PQ 上的 O 点射到 M 上后沿原路反射。线圈通入电流后弹簧长度改变，使 M 发生倾斜，入射光束在 M 上的入射点仍近似处于 PQ 的圆心，通过读取反射光射到 PQ 上的位置，可以测得电流的大小。已知弹簧的劲度系数为 $k$，磁场磁感应强度大小为 $B$，线圈 C 的匝数为 $N$，沿水平方向的长度为 $l$，细杆 D 的长度为 $d$，圆弧 PQ 的半径为 $r$，$r\gg d$，$d$ 远大于弹簧长度改变量的绝对值。
\begin{enumerate}
	\item
	若在线圈中通入的微小电流为 $I$，求平衡后弹簧长度改变量的绝对值 $\Delta x$ 及 PQ 上反射光点与 O 点间的弧长 $s$；
	\item
	某同学用此装置测一微小电流，测量前未调零，将电流通入线圈后，PQ 上反射光点出现在 O 点上方，与 O 点间的弧长为 $s_{1}$，保持其它条件不变，只将该电流反向接入，则反射光点出现在 O 点下方，与 O 点间的弧长为 $s_{2}$。求待测电流的大小。
\end{enumerate}

\onepicture[align=right, w=5cm]{figs/12.png}

%% answer: 
\jdanswer{
\begin{enumerate}
	\item
	$\frac{NBIl}{k}$，$\frac{2NBIlr}{dk}$

	\item
	$\frac{dk(s_{1}+s_{2})}{4NBlr}$
\end{enumerate}
}
%% memo:
\memoanswer{
（1）由题意当线圈中通入微小电流 $I$ 时，线圈中的安培力为 $F=NBIl$

根据胡克定律有

$F=NBIl=k|\Delta x|$

$|\Delta x|=\frac{NBIl}{k}$

如图所示

\onepicture[w=5cm]{figs/12_an.jpg}

设此时细杆转过的弧度为 $\theta$，则可知反射光线转过的弧度为 $2\theta$，又因为

$d\gg\Delta x$，$r\gg d$

则

$\sin\theta\approx\theta$，$\sin2\theta\approx2\theta$

所以有

$\Delta x=d\cdot\theta$

$s=r\cdot2\theta$

联立可得

$s=\frac{2r}{d}\Delta x=\frac{2NBIlr}{dk}$

（2）因为测量前未调零，设没有通电流时偏移的弧长为 $s'$，当初始时反射光点在 O 点上方，通电流 $I'$ 后根据前面的结论可知有

$s_{1}=\frac{2NBI'lr}{dk}+s'$

当电流反向后有

$s_{2}=\frac{2NBI'lr}{dk}-s'$

联立可得

$I'=\frac{dk(s_{1}+s_{2})}{4NBlr}$

同理可得初始时反射光点在 O 点下方结果也相同，故待测电流的大小为

$I'=\frac{dk(s_{1}+s_{2})}{4NBlr}$
}

\item
%% number: 13
%% paperName: 2022年全国高考甲卷第 13 题 10 分
%% typeId: 6
%% type: 计算题
%% chapter: 气体性质
%% point: 气体连接体
%% method: 无
%% score: 10
%% degree: 450
%% duplicateId: 0
%% body:
（1）一定量的理想气体从状态 a 变化到状态 b，其过程如 $p-T$ 图上从 a 到 b 的线段所示。在此过程中\xzanswer{BCE}

\onepicture[w=4.5cm]{figs/13a.png}

\fivechoices[answer=BCE]
{气体一直对外做功}
{气体的内能一直增加}
{气体一直从外界吸热}
{气体吸收的热量等于其对外做的功}
{气体吸收的热量等于其内能的增加量}

（2）如图，容积均为 $V_{0}$、缸壁可导热的 A、B 两汽缸放置在压强为 $p_{0}$、温度为 $T_{0}$ 的环境中；两汽缸的底部通过细管连通，A 汽缸的顶部通过开口 C 与外界相通：汽缸内的两活塞将缸内气体分成 I、Ⅱ、Ⅲ、Ⅳ 四部分，其中第 II、Ⅲ 部分的体积分别为 $\frac{1}{8}V_{0}$ 和 $\frac{1}{4}V_{0}$、环境压强保持不变，不计活塞的质量和体积，忽略摩擦。
\begin{enumerate}
	\item
	将环境温度缓慢升高，求 B 汽缸中的活塞刚到达汽缸底部时的温度；
	\item
	将环境温度缓慢改变至 $2T_{0}$，然后用气泵从开口 C 向汽缸内缓慢注入气体，求 A 汽缸中的活塞到达汽缸底部后，B 汽缸内第 Ⅳ 部分气体的压强。
\end{enumerate}

\onepicture[align=right, w=4.5cm]{figs/13b.png}

%% answer: 
\jdanswer{
\begin{enumerate}
	\item
	$T=\frac{4}{3}T_{0}$

	\item
	$p=\frac{9}{4}p_{0}$
\end{enumerate}
}
%% memo:
\memoanswer{
（1）A．因从 a 到 b 的 $p-T$ 图像过原点，由 $\frac{pV}{T}=C$ 可知从 a 到 b 气体的体积不变，则从 a 到 b 气体不对外做功，选项 A 错误；

B．因从 a 到 b 气体温度升高，可知气体内能增加，选项 B 正确；

CDE．因 $W=0$，$\Delta U>0$，根据热力学第一定律 $\Delta U=W+Q$ 可知，气体一直从外界吸热，且气体吸收的热量等于内能增加量，选项 CE 正确，D 错误。

故选 BCE。

（2）因两活塞的质量不计，则当环境温度升高时，Ⅳ 内的气体压强总等于大气压强，则该气体进行等压变化，则当 B 中的活塞刚到达汽缸底部时，由盖吕萨克定律可得

$\frac{\frac{3}{4}V_{0}}{T_{0}}=\frac{V_{0}}{T}$

解得：$T=\frac{4}{3}T_{0}$

设当 A 中的活塞到达汽缸底部时 Ⅲ 中气体的压强为 $p$，则此时 Ⅳ 内的气体压强也等于 $p$，设此时 Ⅳ 内的气体的体积为 $V$，则 Ⅱ、Ⅲ 两部分气体被压缩的体积为 $V_{0}-V$，则对气体 Ⅳ

$\frac{p_{0}\cdot\frac{3V_{0}}{4}}{T_{0}}=\frac{pV}{2T_{0}}$

对 Ⅱ、Ⅲ 两部分气体

$\frac{p_{0}(\frac{V_{0}}{8}+\frac{V_{0}}{4})}{T_{0}}=\frac{p(V_{0}-V)}{2T_{0}}$

联立解得：$V=\frac{2}{3}V_{0}$，$p=\frac{9}{4}p_{0}$
}

\item
%% number: 14
%% paperName: 2022年全国高考甲卷第 14 题 10 分
%% typeId: 6
%% type: 计算题
%% chapter: 光学
%% point: 全反射
%% method: 无
%% score: 10
%% degree: 550
%% duplicateId: 0
%% body:
（1）一平面简谐横波以速度 $v=2\Ums$ 沿 $x$ 轴正方向传播，$t=0$ 时刻的波形图如图所示，介质中平衡位置在坐标原点的质点 A 在 $t=0$ 时刻的位移 $y=\sqrt 2\Ucm$，该波的波长为\tkanswer[1.6cm]{4}$\Um$，频率为\tkanswer[1.6cm]{0.5}$\UHz$，$t=2\Us$ 时刻，质点 A\tkanswer[2cm]{向下运动}（填“向上运动”、“速度为零”或“向下运动”）。

\onepicture[w=5cm]{figs/14a.png}

（2）如图，边长为 $a$ 的正方形 ABCD 为一棱镜的横截面，M 为 AB 边的中点。在截面所在的平面，一光线自 M 点射入棱镜，入射角为 $60^{\circ}$，经折射后在 BC 边的 N 点恰好发生全反射，反射光线从 CD 边的 P 点射出棱镜，求棱镜的折射率以及 P、C 两点之间的距离。

\onepicture[align=right, w=5cm]{figs/14b.png}

%% answer: 
\jdanswer{
$n=\frac{\sqrt 7}{2}$，$PC=\frac{\sqrt 3-1}{2}a$
}
%% memo:
\memoanswer{
（1）设波的表达式为

$y=A\sin(\frac{2\pi}{\lambda}x+\varphi)$

由题知 $A=2\Ucm$，波图像过点 $(0,\sqrt 2)$ 和 $(1.5,0)$，代入表达式有

$y=2\sin(\frac{\pi}{2}x+\frac{\pi}{4})\Ucm$

即

$\lambda=4\Um$

由于该波的波速 $v=2\Ums$，则

$f=\frac{v}{\lambda}=\frac{2}{4}\UHz=0.5\UHz$

由于该波的波速 $v=2\Ums$，则

$T=\frac{\lambda}{v}=2\Us$

由于题图为 $t=0$ 时刻的波形图，则 $t=2\Us$ 时刻振动形式和零时刻相同，根据“上坡、下坡”法可知质点 A 向下运动。

（2）光线在 M 点发生折射有

$\sin60^{\circ}=n\sin\theta$

由题知，光线经折射后在 BC 边的 N 点恰好发生全反射，则

$\sin C=\frac{1}{n}$

$C=90^{\circ}-\theta$

联立有：$\tan\theta=\frac{\sqrt{3}}{2}$，$n=\frac{\sqrt 7}{2}$

根据几何关系有

$\tan\theta=\frac{MB}{BN}=\frac{a}{2BN}$

$NC=a-BN=a-\frac{a}{\sqrt{3}}$

再由：

$\tan\theta=\frac{PC}{NC}$

解得：

$PC=\frac{\sqrt 3-1}{2}a$
}

\end{enumerate}

\end{document}
